\documentclass[a4paper]{article}
% Español (México) con XeLaTeX: fontspec para fuentes Unicode, polyglossia
% para la división silábica y los nombres del documento en la variante mexicana.
\usepackage{fontspec}
\usepackage{polyglossia}
\setdefaultlanguage[variant=mexican]{spanish}
% Los nombres chinos van como «pinyin (original)»: xeCJK da a los caracteres
% Han su propia fuente; sin ella XeLaTeX los omite con un aviso.
% FandolSong no cubre algunos caracteres de nombres (U+73FA); WenQuanYi Zen Hei sí.
\usepackage[AutoFallBack=true]{xeCJK}
\setCJKmainfont{FandolSong-Regular.otf}
\setCJKfallbackfamilyfont{\CJKrmdefault}{WenQuanYi Zen Hei}
% polyglossia no traduce los nombres de listings.
\AtBeginDocument{\providecommand{\lstlistingname}{}\renewcommand{\lstlistingname}{Listado}%
  \providecommand{\lstlistlistingname}{}\renewcommand{\lstlistlistingname}{Índice de listados}}
\usepackage{amsmath, amssymb}
\usepackage{graphicx}
\usepackage[margin=2.5cm]{geometry}
\usepackage[most]{tcolorbox}
\usepackage{etoolbox}
\usepackage{listings}
\usepackage{hyperref}
\usepackage{booktabs}
\usepackage{subcaption}
\usepackage{float}

\newtcolorbox{knowledgebox}[1]{enhanced, colback=blue!5!white, colframe=blue!75!black, colbacktitle=blue!75!black, coltitle=white, fonttitle=\bfseries, title={#1}, attach boxed title to top left={yshift=-2mm, xshift=2mm}, boxrule=1pt, sharp corners}
\newtcolorbox{importantbox}[1]{enhanced, colback=yellow!10!white, colframe=yellow!80!black, colbacktitle=yellow!80!black, coltitle=black, fonttitle=\bfseries, title={#1}, sharp corners}
\newtcolorbox{warningbox}[1]{enhanced, colback=red!5!white, colframe=red!75!black, colbacktitle=red!75!black, coltitle=white, fonttitle=\bfseries, title={#1}, sharp corners}

\lstset{language=Python, basicstyle=\ttfamily\small, keywordstyle=\color{blue}, stringstyle=\color{red!60!black}, commentstyle=\color{green!60!black}, breaklines=true, frame=single, numbers=left, numberstyle=\tiny\color{gray}, captionpos=b, extendedchars=true}

\newcommand{\notetitle}{{[}LLM Agents F24{]} Course Introduction — Dawn Song \& Xinyun Chen}
\newcommand{\noteauthors}{Elaboradas a partir de materiales públicos del curso}
\newcommand{\notedate}{9 de septiembre de 2024}
\newcommand{\videochannel}{Berkeley RDI}
\newcommand{\videopublishdate}{9 de septiembre de 2024}
\newcommand{\videoduration}{aproximadamente 60 minutos}
\newcommand{\videourl}{https://rdi.berkeley.edu/llm-agents/f24}
\newcommand{\videocoverpath}{cover.jpg}
\newcommand{\repourl}{https://github.com/hqhq1025/ai-course-notes}


% Footer with repo info
\usepackage{fancyhdr}
\pagestyle{fancy}
\fancyhf{}
\fancyfoot[C]{\thepage}
\fancyfoot[R]{\footnotesize\color{gray}\href{\repourl}{AI Course Notes}}
\renewcommand{\headrulewidth}{0pt}
\renewcommand{\footrulewidth}{0pt}

\begin{document}
\begin{titlepage}\centering
{\Large Notas del curso\par}\vspace{1.2cm}
{\huge\bfseries \notetitle\par}\vspace{0.8cm}
{\large \noteauthors\par}\vspace{0.3cm}
{\large \notedate\par}\vspace{1.2cm}
\ifdefempty{\videocoverpath}{{\small Al generar las notas, indique la ruta local de la portada del video.\par}}{\includegraphics[width=0.82\textwidth,height=0.45\textheight,keepaspectratio]{\videocoverpath}\par}
\vfill
\begin{tcolorbox}[width=0.9\textwidth, colback=black!2!white, colframe=black!60, sharp corners]
\textbf{Autor o canal del video}: \videochannel\par
\textbf{Fecha de publicación}: \videopublishdate\par
\textbf{Duración del video}: \videoduration\par
\ifdefempty{\videourl}{}{\textbf{Enlace del video}: \href{\videourl}{\nolinkurl{\videourl}}\par}\par
\textbf{Página del proyecto}: \href{\repourl}{\nolinkurl{\repourl}}\par
\end{tcolorbox}
\end{titlepage}

\tableofcontents\newpage


\section{Introducción al curso: de qué trata esta clase}

Esta clase es la introducción de apertura de «Large Language Model Agents» de Berkeley, en la que Dawn Song presenta el enfoque del curso, el equipo docente y los objetivos de aprendizaje. Ella sitúa a los Agentes en el contexto amplio de la rápida evolución de los LLM: los modelos de lenguaje ya son suficientemente potentes, pero para resolver tareas reales es necesario que el modelo forme un ciclo cerrado con el entorno, las herramientas, la memoria y la retroalimentación.

\begin{figure}[H]
\centering
\includegraphics[page=2,width=0.9\textwidth]{slides.pdf}
\caption{Equipo docente del curso y roles de colaboración. La introducción primero deja claros el cuerpo docente, los ayudantes y los recursos de apoyo del curso.}
\end{figure}

\begin{knowledgebox}{El punto de partida de este curso}
El curso no trata a los Agentes como un conjunto de técnicas de un solo prompt, sino que los define como un sistema completo: el LLM se encarga del razonamiento y la planificación, mientras que los módulos externos se encargan del uso de herramientas, la memoria, la recuperación y la interacción con el entorno. Solo al encadenar estas capacidades puede el LLM pasar de «responder preguntas» a «completar tareas».
\end{knowledgebox}

\subsection{Por qué los modelos grandes avanzan de manera natural hacia los Agentes}

La introducción señala que, en los últimos dos años, la capacidad de los LLM ha mejorado con gran rapidez, pero el modelo de entrada de puro texto a salida de puro texto pronto alcanzó su límite. Las tareas del mundo real por lo general no son preguntas y respuestas de una sola vez, sino que requieren:
\begin{itemize}
    \item interactuar con el entorno, leer páginas, archivos, bases de datos o el estado de dispositivos;
    \item conservar el contexto de la tarea y la memoria a largo plazo;
    \item reintentar tras un fallo, ajustar el plan y corregir el comportamiento a partir de la retroalimentación;
    \item invocar herramientas externas para suplir las carencias del modelo en cómputo, recuperación y ejecución.
\end{itemize}

\begin{importantbox}{El salto clave de LLM a Agente}
Una vez que el modelo cuenta con memoria, uso de herramientas, recuperación y un ciclo de retroalimentación, su unidad de comportamiento deja de ser «responder una frase» para convertirse en «hacer avanzar una tarea a lo largo de una interacción de varios pasos». Esta es también la línea divisoria fundamental entre la investigación de Agentes y las aplicaciones comunes de chatbot.
\end{importantbox}

\subsection{Resumen de la sección}

La introducción del curso deja claro el planteamiento desde el inicio: un Agente no consiste en ponerle otra capa de interfaz de usuario a un LLM, sino en colocar el modelo dentro de un marco de ejecución que observa, planifica, actúa y corrige de manera continua.
\section{Qué es un LLM Agent}

\subsection{Composición del sistema de un Agent}

\begin{figure}[H]
\centering
\includegraphics[page=4,width=0.9\textwidth]{slides.pdf}
\caption{El curso resume los componentes estándar de un Agent con Memory, Tool Use, Retrieval, Reasoning/Planning y Environment Feedback.}
\end{figure}

El panorama que presenta el curso es claro: el centro de un Agent sigue siendo el LLM, pero el LLM ya no funciona de manera aislada, sino que se acopla con sistemas externos mediante varias interfaces.

\begin{itemize}
    \item \textbf{Memory}: conserva el contexto de la tarea, las preferencias del usuario y las trayectorias históricas, para evitar que cada turno de conversación empiece desde cero.
    \item \textbf{Tool use}: permite que el modelo invoque búsquedas, ejecución de código, API, bases de datos u operaciones de navegador.
    \item \textbf{Retrieval}: extrae evidencia externa de documentos y bases de conocimiento, en lugar de depender por completo de la memoria almacenada en los parámetros.
    \item \textbf{Reasoning \& Planning}: divide objetivos complejos en varios pasos, en vez de entregar directamente una salida de un solo turno.
    \item \textbf{Feedback / Environment}: permite que el Agent revise sus propios resultados de ejecución, formando un ciclo cerrado de prueba y error.
\end{itemize}

\begin{knowledgebox}{Por qué conservar estos límites entre módulos}
Si la memoria, la recuperación, la invocación de herramientas y la planeación se comprimen todas en un prompt sumamente largo, el sistema será muy difícil de observar, depurar y optimizar. El valor principal de un framework de Agent está precisamente en separar de forma explícita estas capacidades, de modo que cada componente pueda sustituirse, restringirse y evaluarse por separado.
\end{knowledgebox}

\subsection{Tipos de entornos en los que actúa un Agent}

\begin{figure}[H]
\centering
\includegraphics[page=5,width=0.9\textwidth]{slides.pdf}
\caption{Un LLM Agent no existe solo en interfaces de chat, sino que también puede operar en páginas web, desarrollo de software, flujos de trabajo empresariales y robots, entre otros entornos diversos.}
\end{figure}

El curso hace especial énfasis en la diversidad de entornos de un Agent. Puede trabajar tanto en entornos digitales, como el IDE, el navegador o las plataformas de flujo de trabajo, como extenderse a entornos corpóreos, como robots y tareas de percepción multimodal. Cuanto más abierto es el entorno, más necesita el Agent:
\begin{itemize}
    \item tomar decisiones locales con información incompleta;
    \item recuperarse de los fallos, en vez de terminar de inmediato ante un solo error;
    \item coordinar al mismo tiempo operaciones simbólicas, interacción en lenguaje natural y acciones externas.
\end{itemize}

\subsection{Resumen de la sección}

La definición de un LLM Agent no es algo tan simple como ``un LLM capaz de invocar herramientas'', sino un sistema de ejecución de tareas compuesto por el modelo, capacidades externas y un ciclo de retroalimentación.

\section{¿Por qué estudiar el Agent Framework?}

\subsection{De un solo Agent a la colaboración entre múltiples Agents}

\begin{figure}[H]
\centering
\includegraphics[page=6,width=0.9\textwidth]{slides.pdf}
\caption{La introducción describe la colaboración entre múltiples Agents como un «mecanismo de división del trabajo» en tareas complejas, donde distintos módulos se encargan de distintas subtareas.}
\end{figure}

Cuando la duración de la tarea es suficientemente larga y las subtareas son suficientemente distintas entre sí, un solo Agent suele fallar por la saturación del contexto, la mezcla de responsabilidades y la deriva del razonamiento. Por eso el curso coloca la colaboración entre múltiples Agents entre sus temas centrales:
\begin{itemize}
    \item un planner especializado se encarga de descomponer la tarea;
    \item un specialized executor se encarga de operaciones locales como código, recuperación de información y análisis de tablas;
    \item un verifier / reviewer se encarga de revisar la exactitud factual, la consistencia y la seguridad.
\end{itemize}

\begin{importantbox}{El verdadero valor del Framework}
El valor del Agent framework no está en tener varios conjuntos de API, sino en que ofrece una manera estable de organizar tareas complejas: cómo dividir los módulos, cómo transmitir el contexto, cómo hacer que colaboren varios roles y cómo integrar la recuperación de fallos en el flujo de ejecución.
\end{importantbox}

\subsection{Oportunidades de aplicación y presión de evaluación}

\begin{figure}[H]
\centering
\includegraphics[page=8,width=0.9\textwidth]{slides.pdf}
\caption{El curso presenta lado a lado la generación de código, la automatización de flujos de trabajo, los asistentes personales y la robótica, para mostrar que el límite de aplicaciones de los Agents se está expandiendo rápidamente.}
\end{figure}

\begin{figure}[H]
\centering
\includegraphics[page=9,width=0.9\textwidth]{slides.pdf}
\caption{Puntos de referencia como GAIA, WebArena y SWE-Bench constituyen las primeras coordenadas de evaluación de la investigación sobre Agents.}
\end{figure}

El curso presenta las aplicaciones junto con los benchmark desde el principio, porque los sistemas de Agent son fáciles de convertir en un demo, pero difíciles de poner en producción de manera estable. Los retos que enumera la introducción incluyen:
\begin{itemize}
    \item el razonamiento y la planificación son inestables en tareas complejas;
    \item la eficiencia para aprovechar la retroalimentación del entorno es baja en tareas de largo horizonte;
    \item la comprensión multimodal y el world model todavía no son maduros;
    \item la colaboración entre múltiples Agents, la explicación teórica y la seguridad y privacidad todavía no cuentan con un marco sólido.
\end{itemize}

\begin{warningbox}{La advertencia de riesgo más importante de la introducción}
El fallo de un sistema de Agent normalmente no es «responder mal una pregunta», sino un sesgo acumulado a lo largo de una tarea prolongada: un plan equivocado, una llamada incorrecta a una herramienta, una lectura o escritura errónea del entorno externo, que terminan amplificando un problema pequeño hasta convertirlo en una falla sistémica. Por eso desplegar un Agent exige más observabilidad y evaluación que desplegar una aplicación de LLM común.
\end{warningbox}

\subsection{Resumen de la sección}

La razón para estudiar el Agent framework es muy práctica: determina si el sistema puede manejar tareas reales, y también si esos sistemas pueden evaluarse, depurarse y desplegarse de manera estable.
\section{Estructura del curso, tareas y proyectos}

\subsection{Estratificación del contenido del curso}

\begin{figure}[H]
\centering
\includegraphics[page=11,width=0.9\textwidth]{slides.pdf}
\caption{Este curso divide el contenido en cuatro líneas: capacidades del modelo, Agent framework, aplicaciones y seguridad y ética.}
\end{figure}

La estructura del curso se divide en cuatro niveles:
\begin{enumerate}
    \item \textbf{Model core capabilities}: razonamiento, planeación, comprensión multimodal.
    \item \textbf{LLM agent frameworks}: workflow, tool use, RAG, multi-Agent.
    \item \textbf{Applications}: desarrollo de software, automatización de flujos de trabajo, aplicaciones multimodales y empresariales.
    \item \textbf{Safety \& ethics}: riesgos de despliegue, interacción humano-máquina, problemas de privacidad y control.
\end{enumerate}

\begin{knowledgebox}{Intención de diseño de las direcciones del proyecto}
El proyecto del curso no busca que los estudiantes repitan la construcción de un chatbot, sino que fomenta desarrollarlo en torno a cinco tipos de temas: aplicaciones, puntos de referencia, capacidades fundamentales, seguridad, multi-Agent o sistemas descentralizados. Este diseño en sí mismo indica a los estudiantes que la investigación en Agentes no es una dirección única, sino un espacio de investigación conformado en conjunto por capacidades, sistemas y gobernanza.
\end{knowledgebox}

\subsection{Evaluación y organización del proyecto}

La introducción también presenta con claridad la manera en que se estructura la entrega del curso:
\begin{itemize}
    \item tareas de lectura semanales, que sirven para construir un contexto común entre las distintas clases;
    \item un hands-on lab, para poner realmente en marcha los frameworks y las herramientas;
    \item un semester-long project, que fomenta formar equipos de cinco integrantes y desarrollarlo en torno a un problema completo.
\end{itemize}

En cuanto a los requisitos de calidad, el curso hace corresponder de 1 a 4 unit con distintos niveles de profundidad de implementación, lo cual muestra que da más importancia a «convertir el sistema de Agentes en una entidad de ingeniería que se pueda explicar, mostrar y verificar», y no solo a leer artículos o escribir reseñas conceptuales.

\subsection{Resumen de la sección}

La manera en que se organiza este curso es muy congruente con su propio tema: el curso mismo es un campo de entrenamiento de sistemas de Agentes que va de la lectura y la experimentación hasta la implementación del proyecto.

\section{Síntesis y ampliación}

Esta clase introductoria cumplió tres objetivos: definir la estructura del sistema de un Agente de LLM, explicar por qué el framework de Agentes merece estudiarse por separado, y presentar el mapa de capacidades y la hoja de ruta de proyectos de todo el curso. Su mensaje central es muy directo: las capacidades de los modelos grandes ya son suficientemente sólidas, pero para abordar tareas del mundo real el cuello de botella clave pasa de «si el modelo sabe decirlo» a «si el sistema sabe hacerlo».

\subsection{Lecturas adicionales}

\begin{itemize}
    \item Mialon et al., ``GAIA: a benchmark for General AI Assistants,'' 2023--2024.
    \item Zhou et al., ``WebArena: A Realistic Web Environment for Building Autonomous Agents,'' 2024.
    \item Jimenez et al., ``SWE-Bench: Can Language Models Resolve Real-World GitHub Issues?,'' 2024.
    \item Yao et al., ``ReAct: Synergizing Reasoning and Acting in Language Models,'' ICLR 2023.
\end{itemize}


\end{document}
